% !TeX root = ../papers/nc-p3.tex
\section{階数15の有理な同型商と正則性の二変数への降下}
\label{sec:rank-fifteen-rational-quotient}
\begin{theorem}\label{thm:rank-fifteen-rational-quotient}
第一表示の通常な階数15窓で $U=zk^2$, $V=zm^2$ とする。次の量を取る。
\[
\begin{aligned}
 x&=\frac{V(U+3)}{(U-1)^2}, &y&=\frac{V^3}{(U-1)^4},\\
 p&=\frac{Vzkm(U-9)}{(U-1)^3},&
 q&=\frac{(2n^2+1-U)(z-1)(z\lambda^2-1)}{4z\lambda(U-1)^2}.
\end{aligned}
\]
これらは $U=-3$ も含む全通常点で定義される。
二積の生成次元16の開集合で、二窓が添字保持同型である必要十分条件は
$(x,p,q)$ の一致である。一般点では $x$ と $T=q+2p/(3\sqrt3)$ の一致だけでよい。
従ってこの階数層の添字保持同型類の有理関数体は $k(x,T)$ である。
生成次元の仮定を外しても $(x,p,q)$ の一致は同型の十分条件である。
第二表示には倍率反転で対応する第一表示の式を使う。
\end{theorem}
\begin{proof}
$k$ は標数零の代数閉体とする。前々節の規格化テンソルのパラメータ空間を
$U\ne1$, $V\ne0$ まで広げたものを $X$ とする。具体的には $s,t$ が単元、
$r=s^2,w=t^2$, $K=(r+1/r+w+1/w)/4$, $H=(r+1/r-w-1/w)/4$、
$n^4+(1-K^2)n^2+H^2=0$ であり、$U=K^2$, $V=H^2$ である。
$n$ も単元で、十五係数は非零になる。前々節の192変換は $X$ にも作用する。
実際、変換後も全係数は非零で、八量の復元式が $s,t$ の逆元と $n$ を与え、
$b(a+b)\ne0$ と $LM\ne0$ により $U\ne1$, $V\ne0$ を保つ。
$X$ は既約である。$n\pm H/n$ は $K^2-1\pm2H$ の平方根で、互いから $n$ を復元できる。
$s=2$ で二式を $t$ の既約分数にすると、二つの分子は互いに素な平方自由八次式で、
分母は平方である。従って $k(s,t)$ 上で二つの平方類が独立で、四次式は既約である。

八量の斉次表示で $c=a+b$ として
\[
\begin{gathered}
 P_2=2(a^2+b^2+c^2),\quad P_3=4abc,\quad D_3=(b+c)((b+c)^2-9a^2),\\
 h=4LM,\qquad E=16(L^2-M^2)(B^2-A^2)(D^2-C^2).
\end{gathered}
\]
$a=1$ に戻すと $P_2=U+3$, $P_3=U-1$, $h=H$ である。$E$ は
$(2n^2+1-U)(r-1/r)(w-1/w)/4$ になる。
\path{computations/verify_rank_fifteen_rational_quotient.py} は192変換で
$P_2,E$ が不変、$P_3,D_3$ の符号変化の積が $h$ の符号変化であることを五つの二次式の上で検算する。
従って $x=h^2P_2/P_3^2$, $y=h^6/P_3^4$, $p=h^3D_3/P_3^3$, $q=E/P_3^2$ は不変である。
直接展開すると
\[
 p^2=x^3-27y,\qquad q^2=(1-x)^2-4y
\]
となる。

これらの式で定義される $y\ne0$ の曲面の座標環を $R$ とする。
$y$ を消去し $T_{\pm}=q\pm2p/(3\sqrt3)$ と書けば
\[
 T_+T_- = f(x)=(1-x)^2-4x^3/27=-(x-3)^2(4x-3)/27.
\]
これは既約な超曲面で特異点は高々 $(x,T_+,T_-)=(3,0,0)$ なので、Serreの判定から正規である。
$T_+\ne0$ なら $T_-=f(x)/T_+$ から $p,q$ と $y$ が復元できる。

$X$ は $(x,y)$ 平面の $y\ne0$ の上で有限である。実際
\[
 (U+3)^3=(x^3/y)(U-1)^2,\qquad V^3=y(U-1)^4 .
\]
第一式はモニックで $U=1$ での値は64であるから $U-1$ の逆も整である。
$V$ とその逆、次いで $K,H,s,t,n$ と必要な逆元はモニックな方程式で整になる。
一般の $(x,y)$ に対し $U$ は三通り、各々で $V=x(U-1)^2/(U+3)$ である。
符号付きの $(K,H)$ は四通り、$s,t$ は各四通り、$n$ は四通りで、
これらは全て相異なる $X$ の点になる。従って関数体の次数は768である。

$X$ 上の $t\mapsto1/t$ は $p$ を保ち $q$ を反転し、$(s,t)\mapsto(t,s)$ は
$p$ を反転し $q$ を保つ。いずれも $n,x,y$ を固定するので
$k(x,y,p,q)$ の $k(x,y)$ 上の次数は4である。よって上の整域 $R$ が $k[X]$ に入り、
その分数体が192群の固定体になる（作用の忠実性は前々節の192個の相異なる像から従う）。
$k[X]$ は $R$ 上整で、$R$ は正規だから $k[X]^{G}=R$ である。
有限群の不変式は代数閉体上で相異なる有限軌道を分離する。従って例外ファイバーでも
三値の一致は規格化テンソルの192軌道の一致と同値である。前々節の添字保持同型判定と
$T_-=f(x)/T_+$ を合わせて主張を得る。
\end{proof}

\begin{corollary}\label{cor:rank-fifteen-regularity-quotient}
生成次元を問わず、同じ $(x,y)$ を持つ通常な階数15窓の標準Koszul左右AS正則性は一致する。
この二変数の一般ファイバーにある四つの添字保持同型類は、一つの窓、その一位置移動、
およびそれぞれの反対代数で尽くされる。
\end{corollary}
\begin{proof}
$t\mapsto1/t$ は $\lambda\mapsto1/(z\lambda)$ によって同じ根の反対表示を与え、
通常倍率を保つ。次に $s,t$ の交換では八量の $M,A,C$ の符号だけが変わる。
ラベル $(h,j)$ におけるその符号は
$(-1)^{h\cdot j}(-1)^{h\cdot 2}$ である。第一因子の $X_hZ_j$ の転置が初めの符号を、
対角符号行列の共役が後の符号を与える。第一因子の転置は四文字の位置 $0,2$ を交換するので、
四位置の巡回順序を反転する。同じ根の反対表示の語反転とは一位置移動だけ異なる。
この操作は規格化テンソルで直接成立する。

同じ $(x,y)$ での $(p,q)$ は独立な符号を除けば一致する。必要ならまず $t$ の逆転を、
次に $s,t$ の交換を施すと、目標とは $X$ で同じ192軌道に入る。途中の倍率が非通常でも、
最後の交換と192変換はテンソルの転置と独立共役として実行できる。
従って二つの通常窓の関係空間は巡回移動、語反転と基底変更で対応し、左右正則性が一致する。
一般には両符号がいずれも二択で、定理の必要条件により四類は相異なる。
\end{proof}

\begin{corollary}
正則性を検査する $(x,y)$ 平面では、次の四曲線に落ちる全ての通常窓が
標準Koszul左右AS正則である。
\[
\begin{aligned}
 &4y=(x-1)^2,\qquad y=(x-4)^2,\\
 &x^4-x^3+32x^2y-160xy+256y^2+125y=0,\\
 &x^4+x^3+32x^2y+256y^2-27y=0 .
\end{aligned}
\]
\end{corollary}
\begin{proof}
第一式では $q=0$ で、通常条件から $2n^2+1-U=0$、従って分岐上である。
第二式の左辺から右辺を引いて引き戻すと、単元を除いて
$((U-1)^2-V)F_3$ になる。後の二式の左辺は同様に
$(U-V-2)F_1$, $(U+V)F_2$ となる。
コードはこれらの因数分解を分母を払って検算する。分岐の全点正則性と
定理\ref{thm:rank-fifteen-regular-orbits}を適用する。
\end{proof}
